% !TeX program = pdflatex
\documentclass[11pt,a4paper]{article}
\usepackage[T1]{fontenc}
\usepackage[utf8]{inputenc}
\usepackage[polish]{babel}
\usepackage{lmodern}
\usepackage{amsmath,amssymb,amsthm,mathtools}
\usepackage{microtype}
\usepackage{xcolor}
\usepackage{geometry}
\usepackage{enumitem}
\usepackage{fancyhdr}
\usepackage[unicode]{hyperref}
\hypersetup{colorlinks=true,linkcolor=blue!45!black,urlcolor=blue!45!black}
\geometry{margin=2.5cm}
\setlength{\parindent}{0pt}
\setlength{\parskip}{4pt}
\setlength{\emergencystretch}{2em}
\setlist[enumerate]{topsep=4pt,itemsep=2pt}
\setlist[itemize]{topsep=4pt,itemsep=2pt}
\numberwithin{equation}{section}
\theoremstyle{plain}
\newtheorem{twierdzenie}{Twierdzenie}[section]
\newtheorem{lemat}[twierdzenie]{Lemat}
\newtheorem{stwierdzenie}[twierdzenie]{Stwierdzenie}
\newtheorem{wniosek}[twierdzenie]{Wniosek}
\theoremstyle{definition}
\newtheorem{definicja}[twierdzenie]{Definicja}
\newtheorem{przyklad}[twierdzenie]{Przykład}
\theoremstyle{remark}
\newtheorem{uwaga}[twierdzenie]{Uwaga}
\renewcommand{\proofname}{Dowód}
\pagestyle{fancy}
\fancyhf{}
\fancyhead[L]{\small Paradoks Banacha--Tarskiego}
\fancyhead[R]{\small\ifnum\value{section}>0 Sekcja \thesection\else Wprowadzenie\fi}
\fancyfoot[C]{\thepage}
\renewcommand{\headrulewidth}{0.4pt}
\newcommand{\R}{\mathbb R}
\newcommand{\Z}{\mathbb Z}
\newcommand{\F}{\mathbb F}
\newcommand{\id}{\mathrm{id}}
\newcommand{\red}{\operatorname{red}}
\newcommand{\Fix}{\operatorname{Fix}}
\newcommand{\Orb}{\operatorname{Orb}}
\newcommand{\Stab}{\operatorname{Stab}}
\newcommand{\equivG}{\mathrel{\equiv_G}}
\title{Paradoks Banacha--Tarskiego}
\author{}
\date{1 października 2026}
\begin{document}
\maketitle
\begin{abstract}
Dowodzimy, że domkniętą kulę w $\R^3$ można podzielić na skończenie wiele rozłącznych zbiorów, a następnie za pomocą samych obrotów i przesunięć złożyć z nich dwie rozłączne kule, każdą przystającą do wyjściowej. Zaczynamy od definicji grupy wolnej i jej konstrukcji ze słów. Osobno dowodzimy poprawności redukcji słów, paradoksalnego rozkładu grupy wolnej, istnienia wolnej podgrupy obrotów oraz wszystkich przejść od grupy do sfery i od sfery do pełnej kuli. Nie pomijamy ani słowa pustego, ani punktów na osiach obrotów, ani środka kuli. Wskazujemy miejsce użycia aksjomatu wyboru i wyjaśniamy, dlaczego nie powstaje sprzeczność z addytywnością objętości.
\end{abstract}
\tableofcontents

\section{Cel, założenia i mapa dowodu}
Słowo \emph{paradoks} oznacza tutaj twierdzenie sprzeczne z intuicją, a nie sprzeczność logiczną. Pracujemy w zwykłej teorii mnogości ZFC, czyli z aksjomatem wyboru. Kawałek kuli będzie dowolnym podzbiorem jej punktów: nie musi być bryłą, nie musi być spójny i nie musi mieć objętości w sensie miary Lebesgue'a.

Zakładamy znajomość zbiorów, funkcji, liczb rzeczywistych oraz elementarnego rachunku macierzy i wektorów. Przypominamy potrzebne pojęcia grupowe; nie zakładamy znajomości grup wolnych, teorii działań grup ani teorii rozkładów paradoksalnych. Istnienie liczb rzeczywistych i podstawy algebry liniowej nie są budowane od aksjomatów. Teoria miary pojawia się dopiero w komentarzu po dowodzie i nie jest narzędziem samej konstrukcji.

\subsection{Co dokładnie zostanie udowodnione?}
Przez $\|x\|$ oznaczamy długość euklidesową wektora $x\in\R^3$. Niech
\[
 B^3=\{x\in\R^3:\|x\|\le1\},\qquad S^2=\{x\in\R^3:\|x\|=1\}.
\]
To odpowiednio pełna domknięta kula i jej sfera brzegowa. Nie wolno ich utożsamiać.

\begin{twierdzenie}[Banach--Tarski: podwojenie kuli]\label{tw:banach-tarski}
Istnieją skończenie liczne parami rozłączne zbiory $P_1,\ldots,P_N\subset B^3$ oraz ruchy sztywne $f_i(x)=Q_ix+v_i$, gdzie $Q_i\in\mathrm{SO}(3)$, takie że
\[
 B^3=\bigsqcup_{i=1}^{N}P_i,\qquad
 \bigsqcup_{i=1}^{N}f_i(P_i)
 =(B^3+(-3,0,0))\sqcup(B^3+(3,0,0)).
\]
\end{twierdzenie}
W obu zapisach unie są rozłączne. Każdy kawałek jest użyty dokładnie raz; nie rozciągamy go, nie pomniejszamy i nie tworzymy jego dodatkowej kopii. Konstrukcja podana niżej używa co najwyżej $64$ kawałków; nie jest to twierdzenie o najmniejszej możliwej liczbie kawałków.

Grupę $\mathrm{SO}(3)$, czyli grupę obrotów zachowujących orientację i punkt $0$, zdefiniujemy macierzowo w sekcji~\ref{subsec:so3}. Dodanie $v_i$ pozwala przesuwać kawałki. Odległość między środkami dwóch kul docelowych wynosi $6>2$, więc rzeczywiście są one rozłączne.

\subsection{Siedem etapów}
\begin{enumerate}
\item Zbudujemy grupę wolną $F_2$ z liter $a,a^{-1},b,b^{-1}$.
\item Podzielimy $F_2$ na cztery zbiory, z których dwa odtworzą pierwszą kopię $F_2$, a pozostałe dwa drugą.
\item Pokażemy konkretne dwa obroty, między którymi nie zachodzi żadna nietrywialna relacja słowna. To osadzi $F_2$ w $\mathrm{SO}(3)$.
\item Usuniemy ze sfery przeliczalny zbiór punktów stałych $D$. Na pozostałym zbiorze $X$ przeniesiemy rozkład grupy na każdą orbitę.
\item Wykażemy, że dodanie $D$ nie przeszkadza: sfera i $X$ są równoważne przez rozkład na dwa kawałki.
\item Przedłużymy rozkład sfery wzdłuż promieni, otrzymując rozkład kuli bez środka, a następnie osobno wchłoniemy środek.
\item Umieścimy dwie formalne kopie w dwóch różnych miejscach przestrzeni. Sprawdzimy skończoność, rozłączność i pokrycie.
\end{enumerate}
Najtrudniejszy krok algebraiczny, czyli wolność konkretnych obrotów, ma tu pełny dowód rachunkiem macierzowym modulo $3$. Nie odwołujemy się do twierdzenia o istnieniu „odpowiednich obrotów”.

\section{Grupa wolna: definicja, konstrukcja i przykłady}
\subsection{Minimum języka grupowego}
\begin{definicja}[Grupa]
Grupa to zbiór $G$ z działaniem $(g,h)\mapsto gh$, które jest łączne, ma element neutralny $e$ i dla każdego $g\in G$ ma element odwrotny $g^{-1}$. Oznacza to $(gh)k=g(hk)$, $eg=ge=g$, $gg^{-1}=g^{-1}g=e$. Grupa jest \emph{przemienna}, gdy dodatkowo $gh=hg$ dla wszystkich $g,h$. W tym skrypcie przemienności na ogół nie zakładamy.
\end{definicja}
\begin{przyklad}
$\Z$ z dodawaniem jest grupą: elementem neutralnym jest $0$, a odwrotnością $n$ jest $-n$. Grupa bijekcji danego zbioru ma jako działanie składanie funkcji. W tym drugim przykładzie kolejność może mieć znaczenie. Dla funkcji przyjmujemy $(fg)(x)=f(g(x))$: najpierw działa czynnik prawy.
\end{przyklad}
\begin{definicja}[Homomorfizm, podgrupa, generatory]
Homomorfizm $\varphi:G\to H$ spełnia $\varphi(gh)=\varphi(g)\varphi(h)$. Jeżeli jest różnowartościowy, nazywamy go \emph{osadzeniem}. Podgrupa $K\le G$ to podzbiór zawierający $e$ i zamknięty na iloczyny i odwrotności. Mówimy, że $a,b$ generują $G$, gdy każdy element $G$ jest skończonym iloczynem elementów $a,a^{-1},b,b^{-1}$; dopuszczamy iloczyn pusty, równy $e$.
\end{definicja}
Homomorfizm zachowuje element neutralny i odwrotności. Na przykład $\varphi(e)=\varphi(e)^2$; pomnożenie przez $\varphi(e)^{-1}$ daje $\varphi(e)=e_H$. Następnie $\varphi(g)\varphi(g^{-1})=\varphi(e)=e_H$.

\subsection{Definicja grupy wolnej}
\begin{definicja}[Grupa wolna na dwóch generatorach]\label{def:grupa-wolna}
Para elementów $a,b$ w grupie $F$ jest \emph{wolnym układem generatorów}, jeżeli generuje $F$ i spełnia następującą własność: dla każdej grupy $H$ i każdych $u,v\in H$ istnieje dokładnie jeden homomorfizm $\varphi:F\to H$, $\varphi(a)=u$, $\varphi(b)=v$. Taką grupę oznaczamy $F(a,b)$ lub $F_2$. Własność tę nazywa się \emph{własnością uniwersalną}.
\end{definicja}
Znaczenie definicji jest następujące: obrazy $a$ i $b$ można wybrać zupełnie dowolnie, bez konieczności sprawdzania dodatkowych równań. Gdybyśmy narzucili relację $a^2=e$, to obraz $u$ musiałby spełniać $u^2=e_H$. Gdybyśmy narzucili $ab=ba$, to obrazy musiałyby komutować. Grupa wolna nie narzuca takich relacji.

Ta intuicja nie zastępuje konstrukcji. Teraz zbudujemy grupę, sprawdzimy wszystkie aksjomaty i udowodnimy powyższą własność.

\subsection{Litery, słowa i redukcja}
Bierzemy cztery formalne symbole $\mathcal A=\{a,\bar a,b,\bar b\}$. Na razie kreska nie oznacza obliczonej odwrotności: ustalamy tylko parowanie liter $a\leftrightarrow\bar a$ oraz $b\leftrightarrow\bar b$. Dla litery $s$ oznaczamy sparowaną literę przez $s^{-1}$. Później $\bar a$ rzeczywiście będzie odwrotnością elementu $a$.
\begin{definicja}[Słowo zredukowane]
Słowem jest skończony ciąg liter $s_1\cdots s_n$ z $\mathcal A$. Słowo puste oznaczamy $e$; jego długość wynosi zero. Słowo jest \emph{zredukowane}, jeśli nie zawiera sąsiadującej pary $ss^{-1}$. Redukcja elementarna usuwa taką parę.
\end{definicja}
\begin{przyklad}
Słowo $a\bar a b\bar b ba$ redukuje się następująco: $a\bar a b\bar b ba\longrightarrow b\bar b ba\longrightarrow ba$. Słowo $ab\bar a\bar b$ jest już zredukowane. Nie wolno skasować $a$ z $\bar a$, ponieważ nie sąsiadują. W grupie wolnej nie wolno przestawić liter, aby ułatwić sobie skracanie.
\end{przyklad}

\subsection{Dlaczego wynik redukcji jest jednoznaczny?}
To istotny punkt: bez jednoznaczności nie mielibyśmy dobrze określonego działania. Nie wystarcza powiedzieć „skracamy, ile się da”.
\begin{definicja}[Algorytm stosu]
Czytamy słowo od lewej do prawej, przechowując pomocnicze słowo $r$, początkowo puste. Po przeczytaniu litery $s$:
\begin{itemize}
\item jeśli ostatnią literą $r$ jest $s^{-1}$, usuwamy ją;
\item w przeciwnym razie dopisujemy $s$ na końcu $r$.
\end{itemize}
Wynik po przeczytaniu całego słowa $w$ oznaczamy $\red(w)$. Pomocnicze słowo jest zredukowane po każdym kroku.
\end{definicja}
\begin{lemat}[Odczytanie pary odwrotnej niczego nie zmienia]\label{lem:stos}
Jeśli przed odczytaniem liter $s,s^{-1}$ na stosie jest słowo zredukowane $r$, to po odczytaniu obu liter na stosie znów jest $r$.
\end{lemat}
\begin{proof}
Rozważamy dwa przypadki. Jeśli $r$ nie kończy się literą $s^{-1}$, pierwsza operacja dopisuje $s$, a druga natychmiast usuwa dopisane $s$. Wynik to $r$. Jeśli $r=us^{-1}$, pierwsza operacja pozostawia $u$. Ponieważ $r$ było zredukowane, $u$ nie kończy się literą $s$. Druga operacja dopisuje zatem $s^{-1}$, a nie kasuje końca $u$. Ponownie otrzymujemy $us^{-1}=r$. Działa to również dla pustego $u$.
\end{proof}
\begin{twierdzenie}[Jednoznaczność redukcji]\label{tw:redukcja}
Każde słowo przez kolejne redukcje elementarne prowadzi do dokładnie jednego słowa zredukowanego, którym jest $\red(w)$.
\end{twierdzenie}
\begin{proof}
\emph{Istnienie.} Każde usunięcie pary zmniejsza długość o $2$, więc ciąg usunięć musi się skończyć. Wynik końcowy jest zredukowany. Algorytm stosu również można zrealizować przez takie usunięcia: redukujemy przeczytany prefiks, dopisujemy następną literę i ewentualnie usuwamy jedną końcową parę. Indukcja po długości prefiksu uzasadnia, że końcowy stos jest osiągalnym zredukowanym wynikiem.

\emph{Niezależność od wyboru par.} Słowa $pss^{-1}q$ i $pq$ mają ten sam wynik algorytmu stosu. Po prefiksie $p$ stos jest ten sam; dwie dodatkowe litery nie zmieniają go na mocy lematu~\ref{lem:stos}; następnie czytamy ten sam sufiks $q$.

Zatem każde usunięcie pary zachowuje $\red(w)$. Jeżeli pewien ciąg usunięć kończy się słowem zredukowanym $z$, to $\red(w)=\red(z)=z$. Ostatnia równość zachodzi, ponieważ algorytm stosu nie skraca słowa, w którym nie ma odwrotnych sąsiadów. Każdy wynik końcowy jest więc taki sam.
\end{proof}
\begin{wniosek}\label{wn:redukcja-iloczyn}
Dla dowolnych słów $u,v$ mamy $\red(\red(u)v)=\red(uv)$ oraz $\red(u\red(v))=\red(uv)$.
\end{wniosek}
\begin{proof}
Redukcje prowadzące od $u$ do $\red(u)$ można wykonać wewnątrz początkowego fragmentu $uv$. Analogicznie redukcje $v$ można wykonać wewnątrz końcowego fragmentu. Twierdzenie~\ref{tw:redukcja} mówi, że ostateczny wynik od tego nie zależy.
\end{proof}

\subsection{Działanie grupowe na słowach}
Niech $F_2$ będzie zbiorem wszystkich słów zredukowanych. Definiujemy $u\cdot v=\red(uv)$, gdzie $uv$ po prawej oznacza zwykłe sklejenie napisów.
\begin{stwierdzenie}\label{st:grupa-slow}
To działanie czyni $F_2$ grupą. Elementem neutralnym jest słowo puste $e$, a odwrotnością $s_1\cdots s_n$ jest $s_n^{-1}\cdots s_1^{-1}$.
\end{stwierdzenie}
\begin{proof}
Zamkniętość wynika z tego, że $\red$ daje słowo zredukowane. Łączność nie jest dodatkowym założeniem: z wniosku~\ref{wn:redukcja-iloczyn} otrzymujemy $(u\cdot v)\cdot w=\red(\red(uv)w)=\red(uvw)$ oraz $u\cdot(v\cdot w)=\red(u\red(vw))=\red(uvw)$. Doklejenie słowa pustego nic nie zmienia, zatem jest ono neutralne.

Słowo $s_n^{-1}\cdots s_1^{-1}$ jest zredukowane: sąsiednia para odwrotna w nim odpowiadałaby sąsiedniej parze odwrotnej w pierwotnym słowie. W iloczynie $(s_1\cdots s_n)(s_n^{-1}\cdots s_1^{-1})$ kasujemy kolejno pary od środka na zewnątrz i otrzymujemy $e$. Tak samo dzieje się w odwrotnej kolejności iloczynu.
\end{proof}
\begin{twierdzenie}[Własność uniwersalna konstrukcji]\label{tw:uniwersalnosc}
Zbudowana grupa $F_2$ jest grupą wolną na $a,b$ w sensie definicji~\ref{def:grupa-wolna}.
\end{twierdzenie}
\begin{proof}
Weźmy dowolną grupę $H$ i $u,v\in H$. Zastąpmy litery $a,\bar a,b,\bar b$ odpowiednio elementami $u,u^{-1},v,v^{-1}$, a słowo puste elementem $e_H$. Dla każdego słowa wykonujemy iloczyn tych elementów w tej samej kolejności. Każda usuwana para $ss^{-1}$ ma iloczyn $e_H$, więc redukcja słowa nie zmienia wartości w $H$. Otrzymujemy zatem funkcję $\varphi$ na słowach zredukowanych i mamy $\varphi(\red(wz))=\varphi(wz)=\varphi(w)\varphi(z)$. To homomorfizm o żądanych obrazach $a,b$. Jednoznaczność wynika z tego, że homomorfizm musi zachowywać iloczyny i odwrotności: jego wartość na każdym słowie jest wymuszona przez obrazy liter.
\end{proof}
\begin{uwaga}[Dlaczego mówimy o tej grupie wolnej?]
Jeśli dwie grupy spełniają własność uniwersalną dla dwóch wskazanych generatorów, otrzymujemy homomorfizmy w obie strony zachowujące generatory. Ich złożenia też zachowują generatory, tak samo jak identyczność. Z jednoznaczności we własności uniwersalnej złożenia są identycznościami. Grupy są więc izomorficzne, czyli mają tę samą strukturę grupową poza nazwami elementów.
\end{uwaga}

\subsection{Przykłady i kontrprzykłady}
\begin{przyklad}[Jeden generator]
Przy alfabecie $\{a,\bar a\}$ każde niepuste słowo zredukowane składa się wyłącznie z $a$ albo wyłącznie z $\bar a$: zmiana rodzaju litery tworzyłaby parę odwrotną. Dostajemy elementy $a^n$, $n\in\Z$, i regułę $a^ma^n=a^{m+n}$. Grupa wolna na jednym generatorze jest więc izomorficzna z $(\Z,+)$. Jest przemienna. Wolność nie oznacza automatycznie nieprzemienności.
\end{przyklad}
\begin{przyklad}[Dwa generatory]
W $F_2$ słowa $ab$ i $ba$ są różnymi słowami zredukowanymi, więc $ab\ne ba$. W szczególności $aba^{-1}b^{-1}\ne e$. Jest to przykład komutatora. Niezerowa długość zredukowanego słowa jest tutaj pełnym dowodem, że nie otrzymaliśmy elementu neutralnego. Na przykład $(ab)^{-1}=b^{-1}a^{-1}$, a nie $a^{-1}b^{-1}$. Iloczyn $ab\cdot b^{-1}a^{-1}$ redukuje się do $e$.
\end{przyklad}
\begin{przyklad}[Grupa $\Z^2$ nie jest wolna na swoich dwóch standardowych generatorach]
Dla $u=(1,0)$ i $v=(0,1)$ zachodzi $u+v=v+u$. Zastąpienie $a$ przez $u$, a $b$ przez $v$ daje homomorfizm $F_2\to\Z^2$, ale komutator przechodzi na $(0,0)$. Homomorfizm nie jest różnowartościowy. W $\Z^2$ pamiętamy sumy wykładników, lecz tracimy kolejność liter; w $F_2$ ta kolejność ma znaczenie.
\end{przyklad}
\begin{przyklad}[Nieprzemienność nie wystarcza]
Permutacje $(12)$ i $(23)$ zbioru $\{1,2,3\}$ nie komutują, ale $(12)^2=e$. Nie mogą być wolnymi generatorami. Podobnie wybranie dwóch niekomutujących obrotów w przestrzeni nie dowodzi jeszcze, że generują $F_2$: trzeba wykluczyć wszystkie niepuste słowa zredukowane, nie tylko jeden komutator.
\end{przyklad}
\begin{stwierdzenie}[Przeliczalność]\label{st:przeliczalnosc-f2}
Grupa $F_2$ jest przeliczalnie nieskończona. Słów zredukowanych o długości $n\ge1$ jest dokładnie $4\cdot3^{n-1}$.
\end{stwierdzenie}
\begin{proof}
Pierwszą literę wybieramy na cztery sposoby. Każda następna może być dowolna poza odwrotnością poprzedniej, więc mamy trzy możliwości. Możemy wypisać słowa według długości, a w ramach jednej długości według ustalonego porządku liter. To daje wyliczenie wszystkich elementów. Jest ich nieskończenie wiele, bo $e,a,a^2,\ldots$ są różne.
\end{proof}

\section{Skończone rozkłady i dwie kopie zbioru}
\subsection{Działanie grupy i dozwolone ruchy}
\begin{definicja}[Działanie grupy]
Grupa $G$ działa na zbiorze $Y$, jeśli każdej parze $(g,x)\in G\times Y$ przyporządkowano $gx\in Y$ tak, że $ex=x$ i $(gh)x=g(hx)$. Każde przekształcenie $x\mapsto gx$ jest bijekcją, z odwrotnością $x\mapsto g^{-1}x$. Dla $C\subseteq Y$ piszemy $gC=\{gx:x\in C\}$.
\end{definicja}
Przykłady to lewe mnożenie grupy na sobie, działanie obrotów na sferze oraz działanie ruchów sztywnych na $\R^3$.
\begin{definicja}[Równoważność przez skończony rozkład]\label{def:rownorozklad}
Niech $E,F\subseteq Y$, gdzie $G$ działa na $Y$. Piszemy $E\equivG F$, jeżeli istnieją skończona partycja $E=E_1\sqcup\cdots\sqcup E_k$ oraz $g_1,\ldots,g_k\in G$ takie, że $F=g_1E_1\sqcup\cdots\sqcup g_kE_k$. Słowo \emph{partycja} oznacza, że części są parami rozłączne i pokrywają cały zbiór. Puste części można usuwać. Gdy grupa ruchów jest ustalona, opuszczamy indeks $G$.
\end{definicja}
Tak określony rozkład jest równoważny bijekcji $f:E\to F$, która na każdej ze skończenie wielu części $E_i$ ma postać $f(x)=g_ix$. Rozłączność obrazów daje różnowartościowość, a pokrycie $F$ --- surjektywność.

\textbf{Przykład (dwie części).} Niech grupa $(\Z,+)$ działa na $\Z$ przez przesunięcia. Zbiór $E=\{0,1\}$ jest równoważny zbiorowi $F=\{0,3\}$: część $\{0\}$ pozostawiamy na miejscu, a część $\{1\}$ przesuwamy o $2$. Obrazy $\{0\}$ i $\{3\}$ są rozłączne i pokrywają $F$. Ten przykład pokazuje, że na różnych częściach mogą działać różne elementy grupy.

\begin{lemat}[Składanie rozkładów]\label{lem:skladanie}
Relacja $\equivG$ jest zwrotna, symetryczna i przechodnia. Złożenie rozkładów na odpowiednio $p$ i $q$ części wymaga co najwyżej $pq$ części.
\end{lemat}
\begin{proof}
Zwrotność daje jedna część i ruch $e$. Dla symetrii bierzemy części $g_iE_i$ zbioru docelowego i ruchy $g_i^{-1}$. Przejdźmy do przechodniości. Niech $f:E\to F$ działa jako $g_i$ na $E_i$, a $h:F\to H$ jako $h_j$ na $F_j$. Dzielimy $E$ na zbiory
\begin{equation}\label{eq:3-1}
 E_{ij}'=E_i\cap g_i^{-1}(F_j).
\end{equation}
Dla każdego $x\in E$ istnieje dokładnie jedno $i$ z $x\in E_i$, a następnie dokładnie jedno $j$ z $g_ix\in F_j$. Zatem zbiory $E_{ij}'$ tworzą partycję. Na $E_{ij}'$ złożenie $h\circ f$ jest pojedynczym ruchem $h_jg_i$. Ponieważ złożenie bijekcji jest bijekcją, obrazy tych części są rozłączne i pokrywają $H$. Liczba par $(i,j)$ wynosi $pq$.
\end{proof}

\subsection{Kopie etykietowane: sposób na uczciwe liczenie}
Zapis $E\sqcup E$ jest nieprecyzyjny, jeżeli oba symbole oznaczają ten sam podzbiór tej samej przestrzeni. Wprowadzamy dlatego $2E=E\times\{1,2\}$. Pary $(x,1)$ i $(x,2)$ oznaczają dwa różne miejsca docelowe. Etykieta nie jest dodatkową współrzędną euklidesową.

W rozkładzie między zbiorami etykietowanymi każda część ma jedną ustaloną etykietę źródłową i jedną docelową; przekształcenie ma postać $(x,i)\mapsto(gx,j)$. Zmiany etykiety nie uznajemy za fizyczny ruch. Jest to księgowość rozkładu; na końcu etykiety zastąpimy rzeczywistymi przesunięciami w $\R^3$.

Dowód lematu~\ref{lem:skladanie} działa bez zmian, gdy w przecięciach \eqref{eq:3-1} uwzględniamy zgodność etykiety pośredniej. Jeżeli $E\equivG F$, możemy wykonać ten sam rozkład osobno w obu kopiach, otrzymując $2E\equivG2F$.
\begin{definicja}[Zbiór paradoksalny]\label{def:paradoksalny}
Zbiór $E$ jest \emph{paradoksalny} dla danego działania grupy, jeśli $E\equivG2E$. Równoważnie, istnieje partycja $E=A_1\sqcup\cdots\sqcup A_m\sqcup B_1\sqcup\cdots\sqcup B_n$ i elementy $g_i,h_j\in G$ takie, że
\[
 E=\bigsqcup_{i=1}^{m}g_iA_i,\qquad E=\bigsqcup_{j=1}^{n}h_jB_j.
\]
Pierwsze równanie odtwarza kopię z etykietą $1$, drugie kopię z etykietą $2$. Źródłowe części nie są kopiowane.
\end{definicja}
Niepusty zbiór paradoksalny musi być nieskończony: dla niepustego zbioru skończonego bijekcja $E\to2E$ wymagałaby $|E|=2|E|$. W dalszej części zbudujemy taki rozkład dla $F_2$, a następnie przeniesiemy go na sferę i kulę.
\begin{wniosek}[Przenoszenie paradoksalności]\label{wn:przenoszenie}
Jeśli $E\equivG F$ i $F\equivG2F$, to $E\equivG2E$.
\end{wniosek}
\begin{proof}
Składamy trzy skończone rozkłady: $E\equivG F\equivG2F\equivG2E$. Ostatnia równoważność jest dwukrotnym zastosowaniem odwrotnego rozkładu $F\equivG E$. Lemat~\ref{lem:skladanie} zapewnia, że po złożeniu nadal mamy skończenie wiele części i pojedynczy dozwolony ruch na każdej.
\end{proof}
\subsection{Wchłonięcie jednego poziomu nieskończonego łańcucha}
\begin{lemat}[Przesunięcie łańcucha]\label{lem:lancuch}
Niech $D\subseteq E\subseteq Y$, a $T$ będzie dozwoloną bijekcją $Y$. Załóżmy, że zbiory $D,T(D),T^2(D),\ldots$ są parami rozłączne i wszystkie zawierają się w $E$. Wówczas $E\equivG E\setminus D$ przy użyciu co najwyżej dwóch części.
\end{lemat}
\begin{proof}
Niech $L=\bigcup_{n=0}^{\infty}T^n(D)$. Rozłączność poziomów daje $T(L)=\bigcup_{n=1}^{\infty}T^n(D)=L\setminus D$. Dzielimy $E$ na dwie części $L$ i $E\setminus L$. Pierwszą przenosimy przez $T$, drugą pozostawiamy bez zmian. Obrazy to odpowiednio $L\setminus D$ i $E\setminus L$; są rozłączne i ich suma wynosi $E\setminus D$.
\end{proof}
\begin{uwaga}
Użyliśmy nieskończonej unii do zdefiniowania jednego kawałka $L$. Nie wykonujemy nieskończenie wielu ruchów na różnych poziomach. Na całym $L$ działa jeden i ten sam ruch $T$. Jest to kluczowe rozróżnienie między złożonością kawałka a liczbą kawałków.
\end{uwaga}

\section{Paradoksalny rozkład grupy wolnej}
\subsection{Cztery rodzaje pierwszej litery}
Dla $s\in\mathcal A$ niech $W(s)$ oznacza zbiór niepustych słów zredukowanych zaczynających się literą $s$. Zatem
\begin{equation}\label{eq:4-1}
 F_2=\{e\}\sqcup W(a)\sqcup W(\bar a)\sqcup W(b)\sqcup W(\bar b).
\end{equation}
Każde słowo niepuste ma dokładnie jedną pierwszą literę, więc rozłączność i pokrycie są tu dosłowne.
\begin{lemat}[Kasowanie pierwszej litery]\label{lem:pierwsza-litera}
Mamy $aW(\bar a)=F_2\setminus W(a)$ oraz $bW(\bar b)=F_2\setminus W(b)$.
\end{lemat}
\begin{proof}
Sprawdzimy obie inkluzje pierwszej równości. Słowo z $W(\bar a)$ ma postać $\bar au$, gdzie $u$ jest puste albo nie zaczyna się literą $a$; w przeciwnym razie nie byłoby zredukowane. Mnożenie przez $a$ z lewej kasuje parę $a\bar a$ i daje właśnie $u$. Wynik nie należy do $W(a)$.

Odwrotnie, jeśli $u\notin W(a)$, to $u=e$ albo jego pierwsza litera nie jest $a$. Słowo $\bar au$ jest wtedy zredukowane i należy do $W(\bar a)$, a $a(\bar au)=u$. Dotyczy to także $u=e$: jego poprzednikiem jest słowo $\bar a$. Drugą równość dowodzimy tak samo, zastępując $a$ przez $b$.
\end{proof}
Już teraz mamy dwa równania $F_2=W(a)\sqcup aW(\bar a)$ oraz $F_2=W(b)\sqcup bW(\bar b)$. Nie jest to jeszcze pełny rozkład paradoksalny, ponieważ cztery źródłowe zbiory $W(s)$ nie zawierają $e$. Nie wolno zgubić tego jednego elementu.

\subsection{Włączenie słowa pustego bez dodatkowej kopii}
Połóżmy $P=\{a^n:n=0,1,2,\ldots\}$, gdzie $a^0=e$. Zauważmy dokładnie, że
\begin{equation}\label{eq:4-2}
 aP=P\setminus\{e\}=P\cap W(a).
\end{equation}
Pierwsza równość wynika z $a\cdot a^n=a^{n+1}$ i z różności wszystkich nieujemnych potęg. Druga mówi, że spośród tych potęg tylko $a^0$ nie zaczyna się literą $a$.
\begin{twierdzenie}[Cztery kawałki grupy]\label{tw:cztery-kawalki}
Zbiory
\begin{equation}\label{eq:4-3}
 A_1=W(a)\setminus P,\quad A_2=W(\bar a)\cup P,
 \qquad B_1=W(b),\quad B_2=W(\bar b)
\end{equation}
tworzą partycję $F_2$ i spełniają
\begin{equation}\label{eq:4-4}
 F_2=A_1\sqcup aA_2,\qquad F_2=B_1\sqcup bB_2.
\end{equation}
W szczególności działanie $F_2$ na sobie przez lewe mnożenie jest paradoksalne.
\end{twierdzenie}
\begin{proof}
\emph{Partycja źródła.} W partycji \eqref{eq:4-1} przenosimy dodatnie potęgi $a$ ze zbioru $W(a)$ do zbioru $W(\bar a)$ i dołączamy do niego $e$. Nowy zbiór to $A_2$. Pozostała część $W(a)$ to $A_1$. Nic nie zmieniamy w $W(b)$ i $W(\bar b)$. Zbiory pozostają rozłączne, a teraz pokrywają także $e$.

\emph{Pierwsza kopia.} Korzystając z lematu~\ref{lem:pierwsza-litera} i \eqref{eq:4-2}, otrzymujemy $aA_2=aW(\bar a)\sqcup aP=(F_2\setminus W(a))\sqcup(P\setminus\{e\})$. Pierwszy składnik nie zawiera słów zaczynających się od $a$, a drugi składa się właśnie z dodatnich potęg $a$. Ich unia jest rozłączna. Zbiór $A_1=W(a)\setminus P$ jest rozłączny z obydwoma składnikami i uzupełnia je do całego $F_2$.

\emph{Druga kopia.} Lemat~\ref{lem:pierwsza-litera} od razu daje $B_1\sqcup bB_2=W(b)\sqcup(F_2\setminus W(b))=F_2$. W obu kopiach sprawdziliśmy zarówno pokrycie, jak i rozłączność.
\end{proof}
\begin{center}
\begin{tabular}{c|c|c}
 kawałek źródłowy&ruch&etykieta celu\\\hline
 $A_1$&$e$&$1$\\ $A_2$&$a$&$1$\\ $B_1$&$e$&$2$\\ $B_2$&$b$&$2$
\end{tabular}
\end{center}
\begin{przyklad}[Sprawdzenie kilku punktów]
$e,a,a^2$ należą do $A_2$ i po ruchu $a$ trafiają odpowiednio do $a,a^2,a^3$ w pierwszej kopii. Element $\bar a\in A_2$ trafia do $e$ w pierwszej kopii. Słowo $ab\in A_1$ pozostaje $ab$ w tej samej kopii. Element $\bar b\in B_2$ trafia do $e$ w drugiej kopii. Dwa wystąpienia $e$ w celu mają różnych poprzedników i różne etykiety.
\end{przyklad}

\section{Dwa konkretne obroty generujące grupę wolną}
\subsection{Grupa obrotów i ich punkty stałe}\label{subsec:so3}
\begin{definicja}
Macierz $Q\in M_3(\R)$ jest ortogonalna, jeżeli $Q^TQ=I$. Definiujemy $\mathrm{SO}(3)=\{Q\in M_3(\R):Q^TQ=I,\ \det Q=1\}$. Macierze te zachowują długości i orientację. Działają na $\R^3$ przez $x\mapsto Qx$ oraz, przez ograniczenie, na $S^2$.
\end{definicja}
Rzeczywiście $\|Qx\|^2=x^TQ^TQx=\|x\|^2$. Iloczyn dwóch takich macierzy i odwrotność $Q^{-1}=Q^T$ również spełniają warunki definicji. Mamy więc grupę.
\begin{lemat}[Oś nietrywialnego obrotu]\label{lem:os-obrotu}
Jeśli $Q\in\mathrm{SO}(3)$ i $Q\ne I$, to zbiór wektorów ustalanych przez $Q$ jest prostą przez $0$. W szczególności $Q$ ma dokładnie dwa punkty stałe na $S^2$.
\end{lemat}
\begin{proof}
Najpierw wykażemy istnienie niezerowego wektora stałego. Mamy
\[
 \det(Q-I)=\det Q\det(I-Q^{-1})=\det(I-Q^T)
 =\det(I-Q)=-\det(Q-I).
\]
Ostatni znak minus wynika z nieparzystego wymiaru $3$. Zatem $\det(Q-I)=0$, więc układ $(Q-I)x=0$ ma rozwiązanie $x\ne0$.

Gdyby $Q$ ustalało dwa liniowo niezależne wektory, ustalałoby całą rozpiętą przez nie płaszczyznę. Wybierzmy w niej ortonormalną bazę $u,v$ i jednostkowy wektor $w$ prostopadły do tej płaszczyzny. Ortogonalność implikuje $Qw\perp u,v$ i $\|Qw\|=1$, zatem $Qw=\pm w$. W bazie $u,v,w$ wyznacznik wynosi odpowiednio $1$ albo $-1$. Warunek $\det Q=1$ wymusza $Qw=w$, a wtedy $Q=I$. To sprzeczność. Przestrzeń punktów stałych ma więc wymiar dokładnie jeden. Prosta przez $0$ przecina sferę w dwóch punktach antypodycznych.
\end{proof}
\subsection{Wybór generatorów}
Rozważmy
\begin{equation}\label{eq:5-1}
 A=\frac13\begin{pmatrix}3&0&0\\0&1&-2\sqrt2\\0&2\sqrt2&1\end{pmatrix},\qquad
 B=\frac13\begin{pmatrix}1&-2\sqrt2&0\\2\sqrt2&1&0\\0&0&3\end{pmatrix}.
\end{equation}
Niezerowy blok dwuwymiarowy ma postać $\begin{psmallmatrix}c&-s\\s&c\end{psmallmatrix}$, gdzie $c=1/3$, $s=2\sqrt2/3$ i $c^2+s^2=1$. Jego kolumny są ortonormalne i wyznacznik wynosi $1$. Wobec tego $A,B\in\mathrm{SO}(3)$. Macierz $A$ obraca wokół osi $x$, a $B$ wokół osi $z$, o kąt o cosinusie $1/3$.

Z własności uniwersalnej istnieje homomorfizm $\rho:F_2\to\mathrm{SO}(3)$, $\rho(a)=A$, $\rho(b)=B$. Musimy udowodnić, że jest \emph{różnowartościowy}. To znacznie mocniejsze niż $AB\ne BA$.

\subsection{Rachunek modulo 3 i usunięcie pierwiastków}
Dla liczb całkowitych $m\equiv n\pmod3$ oznacza, że $m-n$ jest podzielne przez $3$. Występują tylko trzy reszty, które wygodnie zapisujemy jako $0,1,-1$. Mamy $2\equiv-1$. Redukcja modulo $3$ zachowuje sumy i iloczyny. Dotyczy to również iloczynów macierzy całkowitych, ponieważ ich wpisy powstają przez sumowanie iloczynów wpisów. Reszty tworzą ciało $\F_3$; potrzebujemy jedynie tego, że iloczyn dowolnej liczby czynników $1$ i $-1$ nie jest zerem. Nie będziemy dzielić przez $3$ po redukcji modulo $3$.

Dla litery $s\in\{a,\bar a,b,\bar b\}$ oznaczmy przez $R_s$ odpowiednio $A,A^{-1},B,B^{-1}$. Macierze $3R_s$ mają wpisy z pierwiastkiem $\sqrt2$. Aby dostać macierze całkowite, wprowadzamy pomocniczo $D_0=\operatorname{diag}(1,\sqrt2,1)$ oraz $N_s=3D_0^{-1}R_sD_0$. Bezpośrednie mnożenie daje
\begin{equation}\label{eq:5-2}
 N_a=\begin{pmatrix}3&0&0\\0&1&-2\\0&4&1\end{pmatrix},\qquad
 N_{\bar a}=\begin{pmatrix}3&0&0\\0&1&2\\0&-4&1\end{pmatrix},
\end{equation}
\begin{equation}\label{eq:5-3}
 N_b=\begin{pmatrix}1&-4&0\\2&1&0\\0&0&3\end{pmatrix},\qquad
 N_{\bar b}=\begin{pmatrix}1&4&0\\-2&1&0\\0&0&3\end{pmatrix}.
\end{equation}
Na przykład wpis $(2,3)$ w $D_0^{-1}(3A)D_0$ wynosi $(1/\sqrt2)(-2\sqrt2)\cdot1=-2$, a wpis $(3,2)$ wynosi $1\cdot2\sqrt2\cdot\sqrt2=4$. Macierze odwrotnych obrotów otrzymujemy z $A^{-1}=A^T$, $B^{-1}=B^T$, co wyjaśnia zmianę znaków.
\begin{uwaga}[Ta zmiana współrzędnych nie porusza kawałków]
$D_0$ nie jest obrotem. Używamy go wyłącznie w dowodzie algebraicznym, aby sprawdzić relacje między rzeczywistymi obrotami $A,B$. Żaden kawałek sfery ani kuli nie zostanie przekształcony przez $D_0$. Nie przemycamy w ten sposób skalowania do twierdzenia.
\end{uwaga}

\subsection{Cztery wektory i jedna tabela}
W rachunku modulo $3$ przyjmijmy wektory kolumnowe
\[
 w_a=\begin{pmatrix}0\\1\\1\end{pmatrix},\quad
 w_{\bar a}=\begin{pmatrix}0\\1\\-1\end{pmatrix},\quad
 w_b=\begin{pmatrix}1\\-1\\0\end{pmatrix},\quad
 w_{\bar b}=\begin{pmatrix}1\\1\\0\end{pmatrix}.
\]
Iloczyn $uv^T$ oznacza macierz, której wpis $(i,j)$ to $u_iv_j$. Z \eqref{eq:5-2}--\eqref{eq:5-3} otrzymujemy
\begin{equation}\label{eq:5-4}
 \overline{N_s}=w_sw_s^T,
\end{equation}
gdzie kreska oznacza redukcję wszystkich wpisów modulo $3$. Na przykład $\overline{N_a}=\begin{psmallmatrix}0&0&0\\0&1&1\\0&1&1\end{psmallmatrix}=w_aw_a^T$ i $\overline{N_b}=\begin{psmallmatrix}1&-1&0\\-1&1&0\\0&0&0\end{psmallmatrix}=w_bw_b^T$. Pozostałe dwie równości sprawdzamy identycznie.

Tabela iloczynów skalarnych $w_s^Tw_t$, już modulo $3$, ma postać
\begin{equation}\label{eq:5-5}
\begin{array}{c|rrrr}
 s\backslash t&a&\bar a&b&\bar b\\\hline
 a&-1&0&-1&1\\
 \bar a&0&-1&-1&1\\
 b&-1&-1&-1&0\\
 \bar b&1&1&0&-1
\end{array}
\end{equation}
Na przekątnej jest $1+1=2\equiv-1$. Dla pary $a,\bar a$ mamy $1-1=0$, a dla $b,\bar b$ także $1-1=0$. Przy dwóch różnych osiach tylko jedna współrzędna daje niezerowy wkład, równy $1$ albo $-1$. Otrzymujemy ważną równoważność:
\begin{equation}\label{eq:5-6}
 w_s^Tw_t=0\quad\Longleftrightarrow\quad t=s^{-1} \qquad(\text{w }\F_3).
\end{equation}
\begin{lemat}[Iloczyn liczników słowa zredukowanego]\label{lem:liczniki}
Jeżeli $s_1\cdots s_n$ jest niepustym słowem zredukowanym, to $\overline{N_{s_1}\cdots N_{s_n}}$ nie jest macierzą zerową.
\end{lemat}
\begin{proof}
Dla wektorów $u,v,x,y$ zachodzi $(uv^T)(xy^T)=u(v^Tx)y^T$. Wpis $(i,j)$ po lewej stronie jest równy $\sum_k u_iv_kx_ky_j=u_i(\sum_kv_kx_k)y_j$, co dowodzi wzoru bez żadnej dodatkowej teorii. Stosując go kolejno do \eqref{eq:5-4}, dostajemy
\begin{equation}\label{eq:5-7}
 \overline{N_{s_1}\cdots N_{s_n}}=
 \left(\prod_{j=1}^{n-1}w_{s_j}^Tw_{s_{j+1}}\right)w_{s_1}w_{s_n}^T.
\end{equation}
Dla $n=1$ iloczyn skalarnych czynników jest pusty i z definicji równy $1$. Ponieważ słowo jest zredukowane, $s_{j+1}\ne s_j^{-1}$. Według \eqref{eq:5-6} każdy czynnik skalarny jest $1$ albo $-1$, więc ich iloczyn również jest niezerowy. Każdy z wektorów $w_s$ ma pewną współrzędną równą $1$. Macierz $w_{s_1}w_{s_n}^T$ ma zatem co najmniej jeden wpis równy $1$. Pomnożenie przez niezerowy skalar nie może wyzerować tego wpisu.
\end{proof}
\begin{twierdzenie}[Wolna podgrupa obrotów]\label{tw:wolne-obroty}
Homomorfizm $\rho:F_2\to\mathrm{SO}(3)$ określony przez $\rho(a)=A$, $\rho(b)=B$ jest różnowartościowy.
\end{twierdzenie}
\begin{proof}
Niech $w=s_1\cdots s_n$ będzie niepustym słowem zredukowanym. Z definicji $N_s$ i kolejnych skróceń $D_0D_0^{-1}$ wynika $N_{s_1}\cdots N_{s_n}=3^nD_0^{-1}\rho(w)D_0$. Gdyby $\rho(w)=I$, prawa strona byłaby równa $3^nI$. Ponieważ $n\ge1$, jej redukcja modulo $3$ byłaby zerowa, wbrew lematowi~\ref{lem:liczniki}. Zatem żadne niepuste słowo zredukowane nie przechodzi na $I$.

Jeśli $\rho(u)=\rho(v)$, to $\rho(v^{-1}u)=I$. Jedynym elementem mającym taki obraz jest $e$, więc $v^{-1}u=e$ i $u=v$. To dokładnie różnowartościowość.
\end{proof}
\begin{uwaga}[Czego rachunek modulo 3 nie twierdzi]
Nie osadziliśmy nieskończonej grupy w skończonej grupie macierzy modulo $3$. Macierze $\overline{N_s}$ są osobliwe. Ponadto $N_sN_{s^{-1}}=9I$, a nie $I$. Badaliśmy podzielność wpisów liczników iloczynu obrotów, a nie homomorfizm grupy obrotów uzyskany przez redukcję ułamka $1/3$.
\end{uwaga}

\section{Od grupy do sfery bez punktów wyjątkowych}
Od tej chwili oznaczamy przez $G=\rho(F_2)\le\mathrm{SO}(3)$ wolną podgrupę z twierdzenia~\ref{tw:wolne-obroty}. Dzięki różnowartościowości możemy utożsamiać elementy $G$ z ich jedynymi słowami zredukowanymi. Litery $a,b$ w działaniu na sferze oznaczają obroty $A,B$.

\subsection{Orbity, stabilizatory i działanie swobodne}
\begin{definicja}
Orbita punktu $x$ to $\Orb(x)=Gx=\{gx:g\in G\}$. Stabilizator punktu to $\Stab(x)=\{g\in G:gx=x\}$. Działanie jest \emph{swobodne}, jeśli $\Stab(x)=\{e\}$ dla każdego punktu rozważanego zbioru.
\end{definicja}
Relacja „należą do tej samej orbity” jest równoważnością: $x=ex$, równanie $y=gx$ daje $x=g^{-1}y$, a z $y=gx$ i $z=hy$ mamy $z=(hg)x$. Orbity są więc parami rozłączne i pokrywają cały zbiór działania.
\begin{uwaga}[Dwa różne znaczenia słowa „wolny”]
Wolna grupa nie musi działać swobodnie! Nasza grupa $G$ jest wolna, ale obrót $A$ ustala dwa punkty osi $x$ na sferze. Aby punkt sfery miał jednoznaczny opis przez słowo, trzeba usunąć wszystkie punkty ustalane przez jakikolwiek nieidentycznościowy element $G$.
\end{uwaga}

\subsection{Przeliczalny zbiór wyjątkowy}
Definiujemy
\begin{equation}\label{eq:6-1}
 D=\bigcup_{g\in G\setminus\{e\}}\Fix_{S^2}(g),\qquad X=S^2\setminus D,
\end{equation}
gdzie $\Fix_{S^2}(g)=\{x\in S^2:gx=x\}$.
\begin{lemat}\label{lem:wyjatkowy}
Zbiór $D$ jest co najwyżej przeliczalny i niezmienniczy względem $G$. Zbiór $X$ jest niezmienniczy, a działanie $G$ na $X$ jest swobodne.
\end{lemat}
\begin{proof}
Grupa $G$ jest przeliczalna na mocy stwierdzenia~\ref{st:przeliczalnosc-f2} i osadzenia $\rho$. Każdy $g\ne e$ jest nieidentycznościowym obrotem i ma dokładnie dwa punkty stałe na sferze, zgodnie z lematem~\ref{lem:os-obrotu}. Przeliczalna unia zbiorów dwuelementowych jest co najwyżej przeliczalna: wyliczamy elementy grupy i dla każdego wypisujemy oba punkty osi. Można najpierw wypisać punkt, którego pierwsza niezerowa współrzędna jest dodatnia, więc ten porządek nie wymaga nowych arbitralnych wyborów.

Jeśli $gx=x$, to dla dowolnego $h\in G$ mamy $(hgh^{-1})(hx)=hx$. Sprzężony element $hgh^{-1}$ nie jest neutralny, gdyż równość $hgh^{-1}=e$ implikowałaby $g=e$. Zatem $hD\subseteq D$. Stosując to samo do $h^{-1}$, dostajemy $D\subseteq hD$, więc $hD=D$. Ponieważ $h$ jest bijekcją sfery, także $hX=X$.

Wreszcie, jeśli $x\in X$ i $gx=x$, to $g\ne e$ oznaczałoby $x\in D$, sprzecznie z definicją $X$. Stabilizator każdego punktu $X$ jest więc trywialny.
\end{proof}

\subsection{Wybór reprezentantów: dokładne miejsce aksjomatu wyboru}
\begin{definicja}[Aksjomat wyboru w potrzebnej postaci]
Dla każdej rodziny niepustych zbiorów istnieje funkcja, która każdemu zbiorowi z tej rodziny przyporządkowuje jeden z jego elementów.
\end{definicja}
Zastosujmy ten aksjomat do rodziny wszystkich orbit działania $G$ na $X$. Wybieramy zbiór $T\subseteq X$ zawierający dokładnie po jednym punkcie z każdej orbity. To jest zasadniczy niekonstruktywny wybór w dowodzie. Nie podajemy algorytmu wyznaczającego $T$.
\begin{lemat}[Jednoznaczne współrzędne orbitowe]\label{lem:wspolrzedne-orbitowe}
Każdy punkt $x\in X$ ma dokładnie jedno przedstawienie $x=gt$, gdzie $g\in G$ i $t\in T$.
\end{lemat}
\begin{proof}
Istnienie wynika z wyboru reprezentanta orbity $x$. Jeśli $gt=hu$, to $t$ i $u$ należą do tej samej orbity, bo $t=g^{-1}hu$. Z wyboru jednego reprezentanta wynika $t=u$. Równość $gt=ht$ daje $(h^{-1}g)t=t$. Swoboda działania wymusza $h^{-1}g=e$, czyli $g=h$.
\end{proof}
\begin{uwaga}
To, że każda orbita jest przeliczalna, nie daje automatycznie przepisu wyboru jednego punktu z każdej orbity. Orbit może być nieprzeliczalnie wiele. Aksjomat wyboru nie jest tu skrótem dla wyliczenia wszystkich punktów sfery.
\end{uwaga}

\subsection{Cztery kawałki na wszystkich orbitach jednocześnie}
Dla $C\subseteq G$ definiujemy $C^T=\{gt:g\in C,\ t\in T\}\subseteq X$. Z jednoznaczności w lemacie~\ref{lem:wspolrzedne-orbitowe} wynika, że rozłączne podzbiory $G$ dają rozłączne zbiory $C^T$. Ponadto
\begin{equation}\label{eq:6-2}
 h(C^T)=(hC)^T\qquad(h\in G),
\end{equation}
ponieważ $h(gt)=(hg)t$. Na wszystkich orbitach działa ten sam obrót $h$, niezależny od wybranego reprezentanta.
\begin{twierdzenie}[Paradoksalność $X$]\label{tw:paradoksalnosc-x}
Zbiór $X=S^2\setminus D$ jest paradoksalny pod działaniem obrotów. Wystarczają cztery części.
\end{twierdzenie}
\begin{proof}
Przenosimy zbiory z \eqref{eq:4-3} na orbity: $C_1=A_1^T$, $C_2=A_2^T$, $C_3=B_1^T$, $C_4=B_2^T$. Tworzą partycję $X$, ponieważ odpowiadające im zbiory grupowe tworzą partycję $G$, a każdy punkt ma jednoznaczne współrzędne $gt$. Z \eqref{eq:4-4} i \eqref{eq:6-2} otrzymujemy $X=C_1\sqcup aC_2$ oraz $X=C_3\sqcup bC_4$. To rozkład $X\equivG2X$. Ruchy to $h_1=I$, $h_2=A$, $h_3=I$, $h_4=B$, a etykiety docelowe to odpowiednio $1,1,2,2$.
\end{proof}
Nie mnożymy liczby kawałków przez liczbę orbit: $C_1$ jest jednym zbiorem zawierającym stosowne punkty wszystkich orbit, i tak samo dla $C_2,C_3,C_4$.

\section{Przywrócenie wszystkich punktów sfery}
Pozostało uzasadnić $S^2\equiv X$. Nie wystarcza zdanie „zbiór przeliczalny jest mały”. Równoważność ma być wykonana skończoną liczbą konkretnych ruchów.

\subsection{Jak uniknąć przeliczalnie wielu zakazów?}
\begin{lemat}\label{lem:nieprzeliczalnosc}
Żaden niepusty otwarty przedział liczb rzeczywistych nie jest przeliczalny. W szczególności można w nim wybrać liczbę poza dowolnym podanym zbiorem co najwyżej przeliczalnym.
\end{lemat}
\begin{proof}
Niech $x_1,x_2,\ldots$ będzie dowolną listą punktów przedziału. Wybierzmy w nim niejednopunktowy przedział domknięty $I_0$. Indukcyjnie wybieramy niejednopunktowy przedział domknięty $I_n\subset I_{n-1}$ o długości co najwyżej połowy poprzedniej, taki że $x_n\notin I_n$. Możemy go umieścić po jednej stronie $x_n$; jeśli $x_n$ nie należy do $I_{n-1}$, dowolny mniejszy podprzedział wystarczy.

Przedziały są zstępujące, domknięte i ich długości dążą do zera. Z zupełności liczb rzeczywistych ich przecięcie zawiera dokładnie jeden punkt $x$: jest nim wspólna granica lewych i prawych końców. Dla każdego $n$ mamy $x\in I_n$, więc $x\ne x_n$. Lista nie obejmowała całego przedziału.
\end{proof}
Sfera też jest nieprzeliczalna: zawiera półokrąg parametryzowany różnowartościowo przez $\theta\mapsto(\cos\theta,\sin\theta,0)$ dla $0<\theta<\pi$. Przeliczalność zbioru trójek indeksów naturalnych wynika z wyliczania ich według sumy indeksów; dla każdej sumy jest tylko skończenie wiele trójek. Dlatego przeliczalna rodzina skończonych zbiorów zakazów ma przeliczalną unię.

\subsection{Obrót rozdzielający przeliczalny zbiór}
\begin{lemat}\label{lem:obrot-rozdzielajacy}
Dla dowolnego co najwyżej przeliczalnego $D\subseteq S^2$ istnieje obrót $R\in\mathrm{SO}(3)$, dla którego zbiory $D,RD,R^2D,\ldots$ są parami rozłączne.
\end{lemat}
\begin{proof}
Jeśli $D=\varnothing$, teza jest oczywista. Niech $D\ne\varnothing$. Wybieramy punkt $u\in S^2$ poza $D\cup(-D)$, co jest możliwe, bo ten ostatni zbiór jest przeliczalny. Żaden punkt $D$ nie leży na osi $\R u$.

Niech $R_u(\theta)$ oznacza obrót wokół tej osi o kąt $\theta$. Można go zdefiniować, wybierając ortonormalną bazę $u,v,w$ o dodatniej orientacji i przyjmując
\[
 R_u(\theta)(\alpha u+\beta v+\gamma w)=\alpha u+
 (\beta\cos\theta-\gamma\sin\theta)v+
 (\beta\sin\theta+\gamma\cos\theta)w.
\]
Ze wzorów dodawania kątów wynika $R_u(\theta)^n=R_u(n\theta)$.

Ustalmy $d,e\in D$ i $n\ge1$. Ponieważ $d$ nie leży na osi, jego składowa w płaszczyźnie $u^\perp$ jest niezerowa. Obrót tej składowej o pełny zakres kątów kreśli okrąg, a punkt na tym okręgu wyznacza kąt jednoznacznie modulo $2\pi$. Zatem albo nie istnieje obrót wokół tej osi przenoszący $d$ na $e$, albo istnieje dokładnie jedna klasa kątów $\alpha$ modulo $2\pi$ o tej własności.

Równanie $R_u(n\theta)d=e$ sprowadza się w drugim przypadku do $n\theta\equiv\alpha\pmod{2\pi}$. W przedziale $0<\theta<2\pi$ ma najwyżej $n$ rozwiązań. Dla każdej trójki $(d,e,n)$ wykluczamy więc skończenie wiele kątów. Trójek jest co najwyżej przeliczalnie wiele. Na mocy lematu~\ref{lem:nieprzeliczalnosc} możemy wybrać $\theta$ poza unią wszystkich wykluczonych kątów.

Dla $R=R_u(\theta)$ mamy teraz
\begin{equation}\label{eq:7-1}
 R^nD\cap D=\varnothing\qquad(n\ge1).
\end{equation}
Jeśli dla $0\le i<j$ istniałby punkt w $R^iD\cap R^jD$, to jego obraz przez $R^{-i}$ należałby do $D\cap R^{j-i}D$, wbrew \eqref{eq:7-1}. Dostajemy żądaną rozłączność wszystkich poziomów.
\end{proof}
\begin{uwaga}
Nowy obrót $R$ nie musi należeć do $G$. To dobrze: $D$ jest $G$-niezmienniczy, więc element $G$ nie mógłby przenieść niepustego $D$ całkowicie poza $D$. Od tego miejsca dla równoważności całej sfery dopuszczamy pełną grupę obrotów $\mathrm{SO}(3)$.
\end{uwaga}
\begin{twierdzenie}[Pełna sfera jest paradoksalna]\label{tw:pelna-sfera}
Sfera $S^2$ jest równoważna przez skończony rozkład dwóm swoim etykietowanym kopiom, przy użyciu samych obrotów wokół $0$. Wystarcza co najwyżej $16$ części.
\end{twierdzenie}
\begin{proof}
Stosujemy lemat~\ref{lem:obrot-rozdzielajacy} do zbioru $D$ z \eqref{eq:6-1}. Niech $L=\bigcup_{n=0}^\infty R^nD$. Wszystkie poziomy leżą na sferze, więc lemat~\ref{lem:lancuch} daje dwuczęściowy rozkład $S^2\equiv X$. Twierdzenie~\ref{tw:paradoksalnosc-x} daje $X\equiv2X$. Przeniesienie paradoksalności z wniosku~\ref{wn:przenoszenie} kończy dowód istnienia rozkładu $S^2\equiv2S^2$.

Poniżej podajemy dokładne części, aby samo odwołanie do przechodniości nie ukrywało konstrukcji ani oszacowania ich liczby. Niech $E_0=L$, $E_1=S^2\setminus L$, $g_0=R$, $g_1=I$. Bijekcja $f:S^2\to X$ ma na $E_i$ postać $g_i$. Zbiory $g_0E_0=L\setminus D$ i $g_1E_1=S^2\setminus L$ tworzą partycję $X$.

Oznaczmy przez $C_j,h_j,\ell(j)$ cztery części, ruchy i etykiety z dowodu twierdzenia~\ref{tw:paradoksalnosc-x}. Dla $i,k\in\{0,1\}$ i $j\in\{1,2,3,4\}$ połóżmy
\begin{equation}\label{eq:7-2}
 H_{ijk}=\{x\in E_i:g_ix\in C_j,\ h_jg_ix\in g_kE_k\}.
\end{equation}
Na tej części działa obrót
\begin{equation}\label{eq:7-3}
 q_{ijk}=g_k^{-1}h_jg_i,
\end{equation}
z etykietą docelową $\ell(j)$. Każde $x$ ma jedyne $i$, następnie jedyne $j$, a potem jedyne $k$. Zatem części \eqref{eq:7-2} tworzą partycję sfery.

Sprawdźmy jeszcze obrazy. Całe odwzorowanie jest złożeniem bijekcji $f$, czteroczęściowej bijekcji $X\to2X$ i odwrotnej bijekcji $f^{-1}$ w każdej kopii. Jest zatem bijekcją $S^2\to2S^2$. Stąd obrazy kawałków są rozłączne wewnątrz każdej etykiety i pokrywają obie sfery. Indeksów jest $2\cdot4\cdot2=16$. Każde $q_{ijk}$ jest złożeniem obrotów, a więc nadal jednym obrotem w $\mathrm{SO}(3)$.
\end{proof}

\section{Od sfery do pełnej kuli}
\subsection{Przedłużenie kawałków wzdłuż promieni}
Niech $B^3_*=B^3\setminus\{0\}$. Dla $H\subseteq S^2$ definiujemy jego stożek promieniowy w kuli bez środka:
\[
 \widehat H=\{rx:x\in H,\ 0<r\le1\}.
\]
Każdy $y\in B^3_*$ ma dokładnie jedno przedstawienie $y=rx$ z $0<r\le1$, $x\in S^2$: mianowicie $r=\|y\|$, $x=y/\|y\|$. Dlatego partycja sfery daje partycję $B^3_*$ przez odpowiednie stożki.
\begin{lemat}[Obroty komutują z przedłużeniem promieniowym]\label{lem:stozki}
Dla $Q\in\mathrm{SO}(3)$ i $H\subseteq S^2$ mamy $Q\widehat H=\widehat{QH}$.
\end{lemat}
\begin{proof}
Liniowość macierzy daje $Q(rx)=rQx$, a ortogonalność daje $\|Qx\|=1$. To dowodzi jednej inkluzji. Jeśli $ry$ należy do $\widehat{QH}$, to $y=Qx$ dla pewnego $x\in H$, więc $ry=Q(rx)\in Q\widehat H$. Mamy także drugą inkluzję.
\end{proof}
\begin{wniosek}\label{wn:kula-bez-srodka}
$B^3_*\equiv2B^3_*$ za pomocą co najwyżej $16$ kawałków i obrotów.
\end{wniosek}
\begin{proof}
Bierzemy stożki nad częściami $H_{ijk}$ z \eqref{eq:7-2} i stosujemy te same obroty $q_{ijk}$. Jednoznaczność kierunku i promienia zapewnia rozłączność źródła, a lemat~\ref{lem:stozki} przenosi rozłączność i pokrycie obrazów z każdej sfery na całą kulę bez środka. Nie zmieniamy promienia żadnego punktu podczas tego etapu.
\end{proof}
\begin{uwaga}[Dlaczego nie dopisać po prostu środka?]
Wszystkie obroty wokół $0$ ustalają $0$. Z jednego takiego punktu nie otrzymamy dwóch punktów oznaczonych różnymi etykietami przez dopisanie go do dwóch części --- przestałyby być rozłączne. Środek wymaga osobnego argumentu i ruchu wokół osi nieprzechodzącej przez $0$.
\end{uwaga}

\subsection{Wchłonięcie środka po małym okręgu}
W płaszczyźnie $z=0$ weźmy okrąg
\[
 \Gamma=\{(1/4,0,0)+(1/4)(\cos t,\sin t,0):t\in\R\}.
\]
Ma środek $c=(1/4,0,0)$ i promień $1/4$. Punkt $0$ należy do $\Gamma$ (dla $t=\pi$). Każdy punkt okręgu ma normę nie większą niż $\|c\|+1/4=1/2$, więc cały okrąg leży wewnątrz $B^3$.

Niech $U$ będzie obrotem wokół osi $z$ o kąt $\alpha=2\pi\sqrt2$. Określmy ruch sztywny
\begin{equation}\label{eq:8-1}
 T_0(x)=c+U(x-c)=Ux+(c-Uc).
\end{equation}
Jest to obrót wokół prostej $c+\R(0,0,1)$; zachowuje okrąg $\Gamma$.
\begin{lemat}\label{lem:wchloniecie-srodka}
Punkty $0,T_0(0),T_0^2(0),\ldots$ są parami różne. W konsekwencji $B^3\equiv B^3_*$ za pomocą dwóch części i ruchów sztywnych zachowujących orientację.
\end{lemat}
\begin{proof}
Gdyby $T_0^m(0)=T_0^n(0)$ dla $m>n$, to $T_0^{m-n}(0)=0$. Punkt $0$ leży na okręgu o dodatnim promieniu, więc wraca na swoje miejsce dokładnie wtedy, gdy kąt obrotu jest całkowitą wielokrotnością $2\pi$. Dostajemy $(m-n)\alpha=2\pi k$, a więc $(m-n)\sqrt2=k$ dla pewnego $k\in\Z$. To sprzeczność z niewymiernością $\sqrt2$.

Dla kompletności: gdyby $\sqrt2=p/q$ w postaci nieskracalnej, to $p^2=2q^2$, więc $p$ byłoby parzyste. Zapis $p=2r$ dawałby $q^2=2r^2$, więc także $q$ byłoby parzyste, wbrew nieskracalności.

Niech $M=\{T_0^n(0):n\ge0\}$. Wszystkie te punkty leżą na $\Gamma\subset B^3$. Stosujemy lemat~\ref{lem:lancuch} dla $D=\{0\}$, $E=B^3$, $T=T_0$. Częściami są $M$ i $B^3\setminus M$, a ruchami $T_0$ i identyczność. Ich obrazy to $M\setminus\{0\}$ oraz $B^3\setminus M$, czyli partycja $B^3_*$.
\end{proof}
\begin{uwaga}
Nie twierdzimy, że $T_0(B^3)=B^3$. Nie jest to potrzebne. Obracamy wyłącznie kawałek $M$, który leży na zachowywanym okręgu. Reszta kuli nie jest poruszana. Dowolny ruch nie musi zachowywać całej kuli, tylko poprawnie przenosić przypisany mu kawałek.
\end{uwaga}

\subsection{Złożenie wszystkich etapów}
\begin{proof}[Dowód twierdzenia~\ref{tw:banach-tarski}]
Lemat~\ref{lem:wchloniecie-srodka} i wniosek~\ref{wn:kula-bez-srodka} dają $B^3\equiv B^3_*\equiv2B^3_*\equiv2B^3$. Przechodniość została udowodniona przez jawne rozdrobnienie części w lemacie~\ref{lem:skladanie}. Nie używamy więc żadnego nieskończonego składania ruchów.

Podajmy również bezpośredni wzór na końcowe części. Zanumerujmy niepuste części rozkładu sfery jako $H_v$, ich obroty jako $q_v$, a etykiety jako $\ell(v)\in\{1,2\}$, gdzie $v$ przebiega zbiór najwyżej $16$ indeksów. Niech $\widehat H_v$ oznacza ich stożki promieniowe. Przyjmijmy $J_0=M$, $J_1=B^3\setminus M$, $t_0=T_0$, $t_1=\id$. Bijekcja $k:B^3\to B^3_*$ działa na $J_s$ jako $t_s$. Zbiory $t_0J_0$ i $t_1J_1$ tworzą partycję $B^3_*$.

Dla $s,t\in\{0,1\}$ i indeksu $v$ definiujemy
\begin{equation}\label{eq:8-2}
 P_{svt}=\{x\in J_s:t_sx\in\widehat H_v,\quad q_vt_sx\in t_tJ_t\}.
\end{equation}
Na tym zbiorze złożenie $k$, paradoksalnego rozkładu $B^3_*$ i odwrotnego $k$ ma postać jednego ruchu
\begin{equation}\label{eq:8-3}
 x\longmapsto t_t^{-1}q_vt_sx,
\end{equation}
z etykietą $\ell(v)$. Dla każdego $x\in B^3$ istnieje jedyne $s$ z $x\in J_s$, jedyne $v$ z $t_sx\in\widehat H_v$ i jedyne $t$ określające, w której części $t_tJ_t$ leży $q_vt_sx$. To dowodzi, że \eqref{eq:8-2} jest partycją całej $B^3$.

Ponieważ każdy etap jest bijekcją między odpowiednimi zbiorami etykietowanymi, obrazy pod ruchami \eqref{eq:8-3} tworzą partycję dwóch formalnych kul. Liczba części jest nie większa niż $2\cdot16\cdot2=64$. Nie ma dodatkowego czynnika $2$ za etykiety: każdy indeks $v$ ma już jedną, określoną etykietę docelową.

Teraz umieszczamy te kopie w przestrzeni. Dla etykiety $1$ dodajemy $v_1=(-3,0,0)$, a dla etykiety $2$ dodajemy $v_2=(3,0,0)$. Ostateczny ruch dla części $P_{svt}$ to $f_{svt}(x)=t_t^{-1}q_vt_sx+v_{\ell(v)}$. Każde $t_s$ jest ruchem $x\mapsto Qx+v$ z $Q\in\mathrm{SO}(3)$, a każde $q_v$ obrotem. Odwrotności i złożenia takich ruchów znów mają tę postać: jeśli $f(x)=Qx+v$, $g(x)=Rx+w$, to $g(f(x))=RQx+(Rv+w)$ oraz $f^{-1}(x)=Q^{-1}x-Q^{-1}v$.

Zatem wszystkie $f_{svt}$ są dozwolonymi ruchami sztywnymi. Ich obrazy wypełniają dokładnie dwie wskazane kule i są parami rozłączne: wewnątrz każdej kuli wynika to z bijekcji, a między kulami z ich rozłącznego położenia. To kończy dowód.
\end{proof}

\subsection{Kula o innym środku i promieniu}
\begin{wniosek}\label{wn:dowolna-kula}
Każdą pełną domkniętą kulę o promieniu $r>0$ można rozłożyć skończenie na dwie rozłączne kule tego samego promienia, używając wyłącznie ruchów sztywnych zachowujących orientację.
\end{wniosek}
\begin{proof}
Dla kuli o środku $c$ użyjmy pomocniczego podobieństwa $S(x)=c+rx$ do zdefiniowania części $S(P_i)$ z dowodu dla kuli jednostkowej. Jeśli ruch jednostkowy miał postać $f_i(x)=Q_ix+v_i$, to $Sf_iS^{-1}(y)=Q_iy+(c-Q_ic+rv_i)$. To nadal izometria: liniowa część wynosi $Q_i$, nie $rQ_i$. Podobieństwo służy tylko przeniesieniu opisu konstrukcji; rzeczywiste kawałki nowej kuli poruszamy powyższymi izometriami, bez skalowania.
\end{proof}

\section{Dlaczego nie przeczy to objętości?}
\subsection{Miara jest określona tylko na pewnych zbiorach}
Miara Lebesgue'a w $\R^3$ przypisuje zbiorom mierzalnym objętość $\lambda(E)\in[0,\infty]$. Do poniższego argumentu potrzebne są tylko jej standardowe własności:
\begin{itemize}
\item kule są mierzalne i $0<\lambda(B^3)<\infty$;
\item izometrie zachowują mierzalność i objętość;
\item skończona unia rozłącznych zbiorów mierzalnych ma objętość równą sumie ich objętości.
\end{itemize}
W standardowej normalizacji $\lambda(B^3)=4\pi/3$, ale konkretna liczba nie jest istotna. Twierdzenie o istnieniu i własnościach miary Lebesgue'a jest tu wyłącznie tłem interpretacyjnym; główny dowód w poprzednich sekcjach nie korzystał z teorii miary.
\begin{stwierdzenie}[Nie wszystkie kawałki mają objętość]\label{st:niemierzalnosc}
W rozkładzie z twierdzenia~\ref{tw:banach-tarski} co najmniej jeden kawałek jest niemierzalny w sensie Lebesgue'a.
\end{stwierdzenie}
\begin{proof}
Załóżmy przeciwnie, że wszystkie $P_i$ są mierzalne. Ponieważ tworzą partycję $B^3$, mamy $\lambda(B^3)=\sum_{i=1}^{N}\lambda(P_i)$. Obrazy $f_i(P_i)$ też są mierzalne, mają te same objętości i tworzą partycję dwóch rozłącznych kul przystających do $B^3$. Zatem
\[
 2\lambda(B^3)=\sum_{i=1}^{N}\lambda(f_i(P_i))=
 \sum_{i=1}^{N}\lambda(P_i)=\lambda(B^3).
\]
Jest to niemożliwe dla dodatniej skończonej liczby $\lambda(B^3)$. Sprzeczne było założenie o mierzalności wszystkich części, a nie twierdzenie Banacha--Tarskiego.
\end{proof}
\begin{uwaga}
Sformułowanie „objętość wzrosła dwukrotnie podczas obracania kawałków” jest mylące. Dla poszczególnych kawałków nie istnieje kompletny zestaw objętości, które można by zsumować. Nie mamy prawa zastosować addytywności miary do dowolnych podzbiorów tylko dlatego, że ich unia jest zwykłą kulą.
\end{uwaga}

\subsection{Co jest, a co nie jest konstruktywne?}
Generatory obrotów są podane jawnymi macierzami. Cztery zbiory słów są zdefiniowane jawnie. Wzory na składanie kawałków też są jawne. Jednak zbiór reprezentantów $T$ został wybrany aksjomatem wyboru. Nie otrzymaliśmy algorytmu, który dla dowolnego punktu kuli rozstrzyga, do którego kawałka należy.

Nie twierdzimy też, że twierdzenie Banacha--Tarskiego jest równoważne pełnemu aksjomatowi wyboru ani że każdy fragment użytego wyboru jest logicznie niezbędny. Wykazaliśmy konkretnie: w ZFC istnieją zbiory i ruchy o opisanych własnościach. Taka jest treść tego dowodu.

Kula jako model ciągłego zbioru punktów nie jest fizycznym przedmiotem zbudowanym z atomów. Rozkład nie opisuje procedury możliwej do wykonania nożem ani technologii powielania materii. Kawałki nie muszą mieć powierzchni, wnętrza czy geometrycznie regularnego brzegu.

\section{Kontrola dowodu i najczęstsze pułapki}
\subsection{Łańcuch zależności}
Każda strzałka w poniższym schemacie została w tekście uzasadniona:
\[
\begin{gathered}
\text{słowa zredukowane}\xrightarrow{\text{tw. \ref{tw:redukcja}, stw. \ref{st:grupa-slow}}}F_2
\xrightarrow{\text{tw. \ref{tw:cztery-kawalki}}}2F_2,\\
F_2\xrightarrow{\text{tw. \ref{tw:wolne-obroty}}}G\le\mathrm{SO}(3)
\xrightarrow{\text{lem. \ref{lem:wyjatkowy}, wybór }T}X,\\
X\xrightarrow{\text{tw. \ref{tw:paradoksalnosc-x}}}2X,
\qquad S^2\xrightarrow{\text{lem. \ref{lem:lancuch}, \ref{lem:obrot-rozdzielajacy}}}X,\\
S^2\xrightarrow{\text{tw. \ref{tw:pelna-sfera}}}2S^2,\\
\text{stożki (lem. \ref{lem:stozki}):}\qquad B^3_*\equiv2B^3_*,\\
B^3\xrightarrow{\text{lem. \ref{lem:wchloniecie-srodka}}}B^3_*
\equiv2B^3_*\equiv2B^3
\xrightarrow{\text{przesunięcia}}\text{dwie rozłączne kule w }\R^3.
\end{gathered}
\]
\subsection{Dziesięć miejsc wymagających szczególnej uwagi}
\begin{enumerate}
\item \textbf{Redukcja nie polega na przestawianiu liter.} Usuwamy tylko sąsiadujące odwrotności. Jednoznaczność wyniku została udowodniona, a nie przyjęta intuicyjnie.
\item \textbf{Słowo puste musi należeć do partycji.} Włącza je zbiór $P$ w konstrukcji \eqref{eq:4-3}.
\item \textbf{Dwa niekomutujące obroty nie wystarczają.} Rachunek modulo $3$ wyklucza wszystkie niepuste słowa zredukowane.
\item \textbf{Nie redukujemy ułamka $1/3$ modulo $3$.} Redukujemy całkowite macierze liczników $N_s$.
\item \textbf{Grupa wolna to nie działanie swobodne.} Zbiór $D$ usuwa stabilizatory inne niż $\{e\}$.
\item \textbf{Wybór reprezentantów nie jest automatyczny.} Wskazaliśmy użycie aksjomatu wyboru przed lematem~\ref{lem:wspolrzedne-orbitowe}.
\item \textbf{Przeliczalne wyjątki nie znikają same.} Wchłania je dwuczęściowy rozkład oparty na obrocie spoza $G$.
\item \textbf{Kawałek złożony z nieskończenie wielu poziomów to nadal jeden kawałek.} Na całej unii poziomów $L$ działa jeden obrót $R$.
\item \textbf{Środek kuli wymaga osobnego ruchu.} Obroty wokół $0$ nigdy go nie przemieszczą; użyliśmy obrotu wokół osi przesuniętej w \eqref{eq:8-1}.
\item \textbf{Formalne kopie nie są jeszcze dwiema kulami w przestrzeni.} Dopiero końcowe przesunięcia o $v_1,v_2$ realizują etykiety jako rzeczywiste rozłączne położenia.
\end{enumerate}
\subsection{Krótkie zadania kontrolne z rozwiązaniami}
\begin{enumerate}
\item \textbf{Do której części należy $a^5$?} Należy do $P$, więc do $A_2$, a nie do $A_1$. Po ruchu $a$ daje $a^6$ w pierwszej kopii. Źródłem $e$ w tej kopii jest $\bar a$.
\item \textbf{Dlaczego z $hgt=hku$ wynika $g=k$ i $t=u$?} Najpierw skracamy działanie $h$, otrzymując $gt=ku$. Punkty $t,u$ reprezentują tę samą orbitę, zatem są równe. Wtedy $(k^{-1}g)t=t$ i swoboda działania daje $k^{-1}g=e$. Bez swobody ostatni wniosek nie byłby uzasadniony.
\item \textbf{Gdzie w rachunku macierzowym wykrywa się redukcję?} Dokładnie w zerowych wpisach tabeli \eqref{eq:5-5}. Sąsiednia para $ss^{-1}$ daje czynnik $w_s^Tw_{s^{-1}}=0$. Brak takich par gwarantuje, że wszystkie czynniki skalarne w \eqref{eq:5-7} są niezerowe.
\item \textbf{Dlaczego nie wystarczy policzyć samej liczby punktów?} Bijekcja dowolnych nieskończonych zbiorów nie musi być składana z kawałków poruszanych ruchami sztywnymi. W twierdzeniu istotne są jednocześnie: skończona liczba kawałków, rozłączność i ograniczenie do izometrii. Sam fakt istnienia bijekcji $B^3\to2B^3$ tego nie zapewnia.
\item \textbf{Czy możemy powtórzyć podwojenie?} Tak. Po pierwszym podwojeniu stosujemy ten sam rozkład do każdej kuli osobno i otrzymujemy cztery kopie, potem osiem itd. Każda ustalona liczba powtórzeń nadal daje skończony rozkład, bo mnożymy tylko skończone liczby kawałków. Nie twierdzimy, że graniczna procedura z nieskończenie wielu powtórzeń jest skończonym rozkładem w sensie definicji~\ref{def:rownorozklad}.
\end{enumerate}

\section*{Uwagi bibliograficzne}
\addcontentsline{toc}{section}{Uwagi bibliograficzne}
Przedstawiony tekst jest samodzielnym opracowaniem klasycznego argumentu: grupa wolna, swobodne orbity, wchłanianie przeliczalnych wyjątków i przedłużenie promieniowe. Nie zakłada ogólniejszej postaci twierdzenia o równoważności dowolnych ograniczonych zbiorów o niepustym wnętrzu. Udowodniona tutaj wersja to pełne podwojenie domkniętej kuli. Szerszy kontekst i inne warianty rozkładów omawiają:
\begin{itemize}
\item Stefan Banach, Alfred Tarski, \emph{Sur la décomposition des ensembles de points en parties respectivement congruentes}, Fundamenta Mathematicae \textbf{6} (1924), 244--277, \href{https://doi.org/10.4064/fm-6-1-244-277}{DOI: 10.4064/fm-6-1-244-277}.
\item Stan Wagon, \emph{The Banach--Tarski Paradox}, Cambridge University Press, 1985, \href{https://www.cambridge.org/core/books/the-banach-tarski-paradox/EA35091E69D798AE0FAE5A4E8386DD08}{katalog wydawcy}.
\end{itemize}
\end{document}
